\section{Computed coordinate-permutation extensions}
\label{sec:party-extensions}

This optional appendix records the splitting obstruction for the projective
coordinate-permutation extension, together with finite examples.  It may be
skipped without affecting the fixed-coordinate transversal-group theorem.

\subsection*{The nonabelian factor set}

Let \(\Gamma\) be the fixed-coordinate projective product-unitary symmetry
group, let \(\widetilde\Gamma\) be the group allowing coordinate permutations,
and let \(\Pi\) be the realized permutation group, so that
\(1\to\Gamma\to\widetilde\Gamma\to\Pi\to1\) is exact.
Operationally, splitting means that the realized coordinate permutations can
be represented by product symmetries forming an actual subgroup, rather than
only up to fixed-coordinate symmetries.

The extension has a section-free outer action
\(\Pi\to\operatorname{Out}(\Gamma)\): conjugation by any lift of
\(\pi\in\Pi\) gives an automorphism of \(\Gamma\), and changing the lift
changes that automorphism by an inner automorphism, because elements of the
kernel act by inner automorphisms.  For a normalized set-theoretic section
\(s(1)=1\), write
\[
  s(\pi)s(\rho)
    =f_s(\pi,\rho)s(\pi\rho),\qquad f_s(\pi,\rho)\in\Gamma .
\]
Then \(f_s(1,\pi)=f_s(\pi,1)=1\), and, with \(\alpha_\pi\) denoting
conjugation by \(s(\pi)\), associativity of three chosen lifts gives
\[
 f_s(\pi,\rho)f_s(\pi\rho,\upsilon)
 =\alpha_\pi\!\left(f_s(\rho,\upsilon)\right)
   f_s(\pi,\rho\upsilon).
\]
If another normalized section is \(s'(\pi)=b(\pi)s(\pi)\) with
\(b(1)=1\), then substituting and commuting no kernel factors gives
\[
 f_{s'}(\pi,\rho)
 =b(\pi)\alpha_\pi\!\left(b(\rho)\right)
   f_s(\pi,\rho)b(\pi\rho)^{-1}.
\]
The order of the factors is essential because \(\Gamma\) need not be
abelian.  The factor set is identically one precisely when the section
preserves multiplication, so the extension splits exactly when \(f_s\) can
be changed to the constant identity factor set by such a normalized \(b\).
That criterion is what the computed rows below verify.

\subsection*{Computed rows}

\begin{corollary}[Computed coordinate-permutation splittings]
\label{cor:computed-party-splitting}
The realized projective coordinate-permutation extension splits for both
non-GRS \(q=11\) pencil classes, the \(q=11\) GRS evaluation-set
representatives
\[
\{0,1,2,3,4,5\},\ \{0,1,2,3,4,6\},\
\{0,1,2,3,5,6\},\ \{0,1,2,3,5,9\},
\]
the enhanced-symmetry \(t=-1\) rows over \(\F_{17}\) and \(\F_{31}\),
and the \(H_3\) representatives
\((q,\tau)=(11,4),(19,5),(29,6),(31,13)\).
For every listed non-GRS or \(H_3\) row, permutation parity acts on the fixed
split torus \(T\) by inversion.  Thus, after choosing an input coordinate and
allowing coordinate permutations to compare the corresponding encoder views,
a realized odd permutation enlarges the possible input symplectic blocks from
\(T\) to
\[
 N(T)=\langle T,J\rangle,\qquad
 J=\begin{pmatrix}0&-1\\1&0\end{pmatrix},\qquad
 J^2=-I,\qquad J\diag(a,a^{-1})J^{-1}=\diag(a^{-1},a).
\]
For the fields in the table, all of which have \(q>3\), this is the full
normalizer of \(T\), and \(N(T)/T\cong C_2\).  More precisely, the extension
\(1\to T\to N(T)\to C_2\to1\) inside \(\mathrm{SL}_2(q)\) does not split in
odd characteristic because every representative of the nontrivial coset has
square \(-I\).  After passing instead to the matrix projectivization
\(\mathrm{PSL}_2(q)=\mathrm{SL}_2(q)/\{\pm I\}\), the induced extension
splits.  This quotient is distinct from quotienting Clifford unitaries by
global phase.  If the coordinate permutation fixes the input
of a chosen encoder, these blocks act as an actual logical enlargement;
otherwise the permutation compares two encoder
views of the same six-leg tensor.
\end{corollary}

\begin{proof}
For each row, an exact finite computation reconstructs the fifteen
four-support shortened planes, exhausts all \(720\) coordinate permutations,
and checks the propagated local symplectic blocks against the full CSS
Lagrangian.  After quotienting by the fixed-coordinate group identified in
Theorem~\ref{thm:logical-phase}, it is enough to test input block \(I\)
on the GRS locus and input blocks \(I,J\) off it; this discards no
candidate.  In the order in which the six \(q=11\) rows are listed above,
their permutation images are
\[
 S_5,\quad S_4,\quad V_4,\quad C_5,\quad S_3,\quad D_{12},
\]
where \(D_{12}\) has order \(12\).  The \(q=17\) and \(q=31\)
enhanced rows have permutation images \(S_4\) and \(S_5\), respectively, and
each listed \(H_3\) image is \(S_5\).

The certificate records the complete normalized factor set for each
section.  It also closes a complement of order \(\lvert\Pi\rvert\) whose
projection to \(\Pi\) is bijective, and verifies that the associated
normalized cochain changes the factor set to the identity.  Hence every
listed extension splits.  Conjugating the fixed torus by the stored lifts
gives the sign action; any odd lift therefore induces the nontrivial
class of \(N(T)/T\).  The complement is an involution in the projective
coordinate-permutation extension group; its determinant-one input block is \(J\), whose
square is \(-I\).
\end{proof}

\begin{table}[htbp]
\centering
\scriptsize
\begin{tabular}{@{}L{.25\textwidth}cL{.20\textwidth}cL{.16\textwidth}@{}}
\toprule
Field and representative & GRS & Fixed-coordinate group & Permutation image & Extension \\
\midrule
\(q=11,\ t=2,\ z=1\) & no & \(\F_q^2\rtimes T\) & \(S_5\) & splits; \(T\to N(T)\) \\
\(q=11,\ t=10,\ z=9\) & no & \(\F_q^2\rtimes T\) & \(S_4\) & splits; \(T\to N(T)\) \\
\(q=11,\{0,1,2,3,4,5\}\) & yes & \(\F_q^2\rtimes\mathrm{SL}_2(q)\) & \(V_4\) & splits \\
\(q=11,\{0,1,2,3,4,6\}\) & yes & \(\F_q^2\rtimes\mathrm{SL}_2(q)\) & \(C_5\) & splits \\
\(q=11,\{0,1,2,3,5,6\}\) & yes & \(\F_q^2\rtimes\mathrm{SL}_2(q)\) & \(S_3\) & splits \\
\(q=11,\{0,1,2,3,5,9\}\) & yes & \(\F_q^2\rtimes\mathrm{SL}_2(q)\) & \(D_{12}\) & splits \\
\(q=17,\ t=-1\) & yes & \(\F_q^2\rtimes\mathrm{SL}_2(q)\) & \(S_4\) & splits \\
\(q=31,\ t=-1\) & no & \(\F_q^2\rtimes T\) & \(S_5\) & splits; \(T\to N(T)\) \\
\(H_3:(q,\tau)=(11,4)\) & no & \(\F_q^2\rtimes T\) & \(S_5\) & splits; \(T\to N(T)\) \\
\(H_3:(q,\tau)=(19,5)\) & no & \(\F_q^2\rtimes T\) & \(S_5\) & splits; \(T\to N(T)\) \\
\(H_3:(q,\tau)=(29,6)\) & no & \(\F_q^2\rtimes T\) & \(S_5\) & splits; \(T\to N(T)\) \\
\(H_3:(q,\tau)=(31,13)\) & no & \(\F_q^2\rtimes T\) & \(S_5\) & splits; \(T\to N(T)\) \\
\bottomrule
\end{tabular}
\caption{The twelve certificate-checked coordinate-permutation extensions.
Here \(D_{12}\) has order \(12\), and \(T\to N(T)\) means that odd
coordinate permutations supply the nontrivial coset of \(N(T)/T\) when
encoder views are compared at a chosen input coordinate.}
\label{tab:party-extension-census}
\end{table}

At \(q=11\), the \(t=2,z=1\) Clebsch/\(H_3\) class has permutation image
\(S_5\), while the \(t=10,z=9\) non-GRS class has permutation image \(S_4\).
Both extensions split, and odd coordinate permutations invert \(T\).
After choosing an input coordinate, the possible one-qudit projective groups
across the resulting encoder views therefore form
\(\F_q^2\rtimes N(T)\).  This is the logical group of a fixed encoder
only for the subgroup fixing its input; a general coordinate permutation
compares different encoder views.  The GRS fixed-coordinate logical group and
this non-GRS group have orders \(159720\) and \(2420\), respectively.

Arithmetic symmetry jumps do not change which branch of the group theorem
applies.  The characteristic-\(17\) tetrahedral specialization lies on the GRS
locus and retains the full \(\mathrm{SL}_2(17)\) fixed-coordinate image; its permutation
image is \(S_4\), and the computed extension splits.  The
characteristic-\(31\) \(A_5\) specialization is non-GRS, so its
fixed-coordinate image is the split torus; its \(S_5\) permutation image also
splits, and odd permutations enlarge the possible input blocks to \(N(T)\).
